% !TEX root = ../main.tex
\section{Construction}\label{sec:construction}

We now present the \textsc{zk-eidas} construction in full.
Section~\ref{sec:model} defines the system model, roles, and security properties.
Section~\ref{sec:parsing} describes the credential parsing circuit that ingests raw mdoc CBOR bytes and produces a zero-knowledge proof.
Section~\ref{sec:predicates} formalises the predicate system used for selective disclosure.
Section~\ref{sec:nullifiers} specifies nullifiers and holder binding.

%% ───────────────────────────────────────────────────────────────────
\subsection{System Model and Definitions}\label{sec:model}

\paragraph{Roles.}
The construction involves four parties:
\begin{itemize}
  \item \textbf{Issuer $\mathcal{I}$}. Issues mobile driving licence (mdoc) credentials conforming to ISO/IEC~18013-5~\cite{iso18013}. The issuer signs the Mobile Security Object (MSO) with ECDSA over the NIST~P-256 curve. Our construction requires \emph{no modifications} to the issuer.
  \item \textbf{Holder $\mathcal{H}$}. Possesses an mdoc credential and generates zero-knowledge proofs of predicate satisfaction. The holder runs the proving algorithm on their device.
  \item \textbf{Verifier $\mathcal{V}$}. Receives a proof and verifies it. The verifier learns only the Boolean result of each requested predicate --- no attribute values are revealed.
  \item \textbf{Escrow Authority $\mathcal{E}$}. Holds an ML-KEM-768 decapsulation key~\cite{fips203}. The authority can decrypt the holder's escrowed identity fields only when presented with a valid court order or arbitration ruling.
\end{itemize}

\paragraph{Security properties.}
We require the following properties, stated semi-formally below.

\begin{definition}[Completeness]\label{def:completeness}
  For every honestly issued credential $C$, set of predicates $\{P_i\}$ that $C$ satisfies, and honestly generated escrow ciphertext, the proof produced by an honest holder is accepted by the verifier with probability~$1$.
\end{definition}

\begin{definition}[Soundness]\label{def:soundness}
  For any probabilistic polynomial-time adversary $\mathcal{A}$, the probability that $\mathcal{A}$ produces a proof $\pi$ that is accepted by the verifier for a statement $x \notin \mathcal{L}$ is bounded by a negligible function $\mathsf{negl}(\lambda)$:
  \[
    \Pr\bigl[\mathsf{Verify}(\mathsf{vk}, x, \pi) = 1 \;\wedge\; x \notin \mathcal{L}\bigr] \leq \mathsf{negl}(\lambda).
  \]
  In our setting, $x \notin \mathcal{L}$ means either the credential signature is invalid, or at least one predicate is not satisfied by the actual attribute value.
\end{definition}

\begin{definition}[Zero-Knowledge]\label{def:zk}
  There exists a probabilistic polynomial-time simulator $\mathcal{S}$ such that for every verifier $\mathcal{V}^*$, the view of $\mathcal{V}^*$ when interacting with an honest prover on witness $w$ is computationally indistinguishable from the output of $\mathcal{S}(x)$ given only the statement:
  \[
    \bigl\{\mathsf{View}_{\mathcal{V}^*}(\mathcal{H}(w) \leftrightarrow \mathcal{V}^*)\bigr\}
    \;\approx_c\;
    \bigl\{\mathcal{S}(x)\bigr\}.
  \]
  In particular, $\mathcal{V}^*$ learns nothing about the credential beyond the truth value of the requested predicates.
\end{definition}

\begin{definition}[Unlinkability]\label{def:unlinkability}
  Let $\pi_0$ and $\pi_1$ be two proofs generated by the same holder from the same credential under different verifier scopes (i.e., distinct contract hashes $c_0 \neq c_1$). No probabilistic polynomial-time distinguisher $\mathcal{D}$ can determine whether $\pi_0$ and $\pi_1$ originate from the same credential with advantage greater than $\mathsf{negl}(\lambda)$.
\end{definition}

\begin{definition}[Escrow Correctness]\label{def:escrow-correct}
  If the holder honestly encrypts identity fields $\{f_i\}$ under the authority's ML-KEM-768 public key, then the authority recovers $\{f_i\}$ exactly upon decapsulation and AES-256-GCM decryption.
\end{definition}

\begin{definition}[Non-Frameability]\label{def:non-frame}
  No coalition of a malicious authority $\mathcal{E}^*$ and a malicious verifier $\mathcal{V}^*$ can produce a valid escrow ciphertext that, when decrypted, yields identity fields belonging to an honest holder $\mathcal{H}$ who did not participate in the protocol, except with probability $\mathsf{negl}(\lambda)$.
\end{definition}

\paragraph{Threat model.}
We consider the following adversarial capabilities:
\begin{itemize}
  \item \emph{Verifiers} are honest-but-curious: they follow the protocol but attempt to extract information from transcripts.
  \item \emph{Holders} are potentially malicious: they may attempt to prove false statements or use forged credentials.
  \item \emph{The escrow authority} is honest-but-curious: it holds the decapsulation key faithfully but may attempt to learn holder identities from escrowed data without authorisation.
  \item \emph{Collusion}: even if a verifier and the escrow authority collude, they cannot break unlinkability (Definition~\ref{def:unlinkability}) without a court order that triggers decapsulation. The scoped nullifier mechanism ensures that the authority sees only an encrypted blob, and the verifier sees only a nullifier hash --- neither alone nor together suffices to link presentations across different scopes.
  \item \emph{Issuers} are assumed honest. A malicious or compromised issuer could craft credentials with colliding digests in the MSO or embed subliminal information in the ECDSA signature nonce. These attacks are outside our scope: mitigating them requires issuer-side controls (e.g., deterministic ECDSA nonces per RFC~6979) rather than proving-system changes.
\end{itemize}

%% ───────────────────────────────────────────────────────────────────
\subsection{Credential Parsing Circuit}\label{sec:parsing}

The core contribution of \textsc{zk-eidas} is a \emph{monolithic} Longfellow circuit~\cite{longfellow} that ingests raw mdoc CBOR bytes and performs all verification steps --- CBOR parsing, digest chain validation, ECDSA P-256 signature verification, predicate evaluation, nullifier derivation, and escrow digest computation --- within a single proof.  No trusted preprocessing step transforms the credential before it enters the circuit.

Algorithm~\ref{alg:prove} presents the proving procedure.

\begin{algorithm}[t]
  \caption{$\mathsf{Prove}$: mdoc selective-disclosure proof generation}\label{alg:prove}
  \begin{algorithmic}[1]
    \Require mdoc bytes $m$; issuer public key $\mathit{pk} = (x,y) \in \mathbb{F}_p^2$ (P-256); session transcript $\mathit{st}$; requested attributes $\{(ns_i, id_i, v_i, t_i)\}_{i=1}^{k}$; current time $\mathit{now}$; contract hash $c \in \{0,1\}^{64}$; escrow fields $\{e_j\}_{j=0}^{7} \subset \{0,1\}^{256}$
    \Ensure proof bytes $\pi$; nullifier hash $\eta$; binding hash $\beta$; escrow digest $\delta$
    \Statex
    \State Parse CBOR structure from $m$
      \Comment{in-circuit CBOR decoder}
    \State Extract $\mathsf{IssuerAuth}$ (COSE\_Sign1 envelope) from $m$
    \State Extract MSO $\gets$ payload of $\mathsf{IssuerAuth}$
    \For{each requested attribute $(ns_i, id_i, v_i, t_i)$}
      \State Locate $\mathsf{IssuerSignedItem}$ for $(ns_i, id_i)$ in the CBOR structure
      \State $d_i \gets \mathsf{SHA\textnormal{-}256}(\mathsf{IssuerSignedItem})$
      \State \textbf{assert} $d_i = \mathsf{MSO.valueDigests}[ns_i][\mathsf{digestID}]$
        \Comment{digest chain}
      \State $v'_i \gets$ extract $\mathsf{elementValue}$ from CBOR encoding
      \State Evaluate predicate $(v'_i, v_i, t_i)$
        \Comment{see \S\ref{sec:predicates}}
    \EndFor
    \State Verify $\mathsf{ECDSA}_{pk}$ signature on $\mathsf{IssuerAuth}$
      \Comment{P-256, see below}
    \State $\eta \gets \mathsf{SHA\textnormal{-}256}(m \,\|\, c)$
      \Comment{scoped nullifier}
    \State $\beta \gets \mathsf{SHA\textnormal{-}256}(\mathsf{attr}_0[0..31])$
      \Comment{holder binding}
    \State $\delta \gets \mathsf{SHA\textnormal{-}256}(e_0 \,\|\, e_1 \,\|\, \cdots \,\|\, e_7)$
      \Comment{escrow digest}
    \State \Return $(\pi, \eta, \beta, \delta)$
  \end{algorithmic}
\end{algorithm}

\paragraph{Single-circuit architecture.}
A critical design choice is that all steps 1--14 of Algorithm~\ref{alg:prove} execute inside a \emph{single} Longfellow circuit instance.  The circuit takes raw CBOR bytes as private input and produces $({\eta}, {\beta}, {\delta})$ as public outputs.  The Sumcheck+Ligero proof system~\cite{sumcheck,ligero} then proves knowledge of a satisfying assignment.  This monolithic approach eliminates the need for intermediate commitments or multi-circuit composition, keeping the critical path shallow (circuit depth~22 for the signature sub-circuit, depth~21 for the hash sub-circuit).

\paragraph{Circuit topology.}
Longfellow compiles the circuit into two sub-circuits that are composed internally:
\begin{enumerate}
  \item A \emph{hash sub-circuit} over $\mathsf{GF}(2^{128})$ containing ${\sim}\,4 \times 10^6$ wires and ${\sim}\,99{,}000$ public inputs.  This sub-circuit handles CBOR parsing, SHA-256 digest chain verification, predicate evaluation, nullifier computation, binding hash, and escrow digest.
  \item A \emph{signature sub-circuit} over $\mathsf{GF}(2^{128})$ containing ${\sim}\,224{,}000$ wires and ${\sim}\,3{,}700$ public inputs.  This sub-circuit handles ECDSA~P-256 signature verification.
\end{enumerate}
The two sub-circuits are proven independently (each via Sumcheck+Ligero) and their proofs are concatenated.  Shared values between sub-circuits are committed to in both proofs.

\paragraph{ECDSA in-circuit.}
P-256 signature verification is the most expensive component, requiring ${\sim}\,17\,\text{ms}$ of the ${\sim}\,280\,\text{ms}$ proving time for the signature sub-circuit.  Working over $\mathsf{GF}(2^{128})$ provides efficient native SHA-256 gates, but ECDSA requires field arithmetic over $\mathbb{F}_p$ where $p$ is the P-256 prime.  We handle this via \emph{range-checked limb decomposition}: each 256-bit P-256 field element is decomposed into limbs that fit within the $\mathsf{GF}(2^{128})$ arithmetic, with range checks ensuring no overflow.  Scalar multiplication is performed using a double-and-add circuit with constant-time conditional moves.

%% ───────────────────────────────────────────────────────────────────
\subsection{Predicate System}\label{sec:predicates}

Each attribute carries a \emph{verification type} $t \in \{0,1,2,3\}$ that determines which predicate the circuit enforces on the attribute value.  Table~\ref{tab:predicates} lists the four primitive types.

\begin{table}[t]
  \centering
  \caption{Primitive predicate types.  Here $v$ is the actual attribute value extracted from the credential, and $v'$ is the expected value supplied by the verifier.}\label{tab:predicates}
  \begin{tabular}{clll}
    \toprule
    Code & Type & Constraint & Semantics \\
    \midrule
    0 & Equality      & $v = v'$    & Value matches exactly \\
    1 & Less-or-equal & $v \leq v'$ & Value at most $v'$ \\
    2 & Greater-or-equal & $v \geq v'$ & Value at least $v'$ \\
    3 & Not-equal     & $v \neq v'$ & Value differs from $v'$ \\
    \bottomrule
  \end{tabular}
\end{table}

\paragraph{In-circuit comparison.}
Algorithm~\ref{alg:predicate} describes how the circuit evaluates a predicate.  All comparisons are \emph{lexicographic} over the byte-level representation, using a binary comparator sub-circuit.

\begin{algorithm}[t]
  \caption{$\mathsf{EvalPredicate}$: in-circuit predicate evaluation}\label{alg:predicate}
  \begin{algorithmic}[1]
    \Require Actual value $v \in \{0,1\}^n$; expected value $v' \in \{0,1\}^n$; type $t \in \{0,1,2,3\}$
    \Ensure Assertion that the predicate holds
    \Statex
    \State $\mathit{is\_leq} \gets \mathsf{CMP.leq}(n,\, v,\, v')$
      \Comment{$1$ iff $v \leq v'$ lexicographically}
    \State $\mathit{is\_geq} \gets \mathsf{CMP.leq}(n,\, v',\, v)$
      \Comment{$1$ iff $v' \leq v$, i.e.\ $v \geq v'$}
    \State $\mathit{is\_eq} \gets \mathit{is\_leq} \wedge \mathit{is\_geq}$
      \Comment{$v = v'$ iff both hold}
    \State $\mathit{pass} \gets \mathsf{MUX}\bigl(t,\; \mathit{is\_eq},\; \mathit{is\_leq},\; \mathit{is\_geq},\; \neg\,\mathit{is\_eq}\bigr)$
    \State \textbf{assert} $\mathit{pass} = 1$
  \end{algorithmic}
\end{algorithm}

The multiplexer $\mathsf{MUX}$ selects the appropriate Boolean based on the type code~$t$: equality ($t=0$) checks $\mathit{is\_eq}$; less-or-equal ($t=1$) checks $\mathit{is\_leq}$; greater-or-equal ($t=2$) checks $\mathit{is\_geq}$; not-equal ($t=3$) checks $\neg\,\mathit{is\_eq}$.  The assertion is enforced as a hard constraint --- if it fails, the circuit is unsatisfiable and no valid proof can be produced.

\paragraph{Composed predicates.}
While the circuit supports only the four primitive types, the API layer composes them into richer predicates before circuit entry:
\begin{itemize}
  \item \textbf{Range} $[\ell, h]$: decomposed into two attribute requests on the same claim --- one with type $\mathsf{GEQ}(\ell)$ and one with type $\mathsf{LEQ}(h)$.  Both must pass for the proof to succeed.
  \item \textbf{Set membership} $\{v_1, \ldots, v_k\}$: decomposed into $k$ attribute requests with type $\mathsf{EQ}$, combined with a logical~OR at the API layer (at least one must match).  The cost is linear in~$k$: each element adds one attribute slot to the circuit.  For small sets (e.g., EU nationality among 27~Member States) this is practical; for large sets, a Merkle-based membership proof would amortise the cost to $O(\log k)$ and is a natural extension.
  \item \textbf{Boolean combinators}: arbitrary $\mathsf{AND}$ and $\mathsf{OR}$ trees over sub-predicates, resolved by the facade before generating the attribute list passed to the circuit.
\end{itemize}

\paragraph{Date handling.}
CBOR-encoded date attributes in mdoc credentials use the ISO~8601 \emph{full-date} profile (\texttt{YYYY-MM-DD}), e.g., \texttt{"1998-09-04"}.  A key observation is that strings in this fixed-width format sort correctly under lexicographic comparison: for dates $d_1 < d_2$ in calendar order, the byte representation satisfies $d_1 <_{\mathrm{lex}} d_2$.  (This property does \emph{not} hold for arbitrary ISO~8601 representations that include time-zone offsets.)  This allows the circuit to perform date comparisons (e.g., ``birth date $\leq$ cutoff date'' for age proofs) using the same $\mathsf{CMP.leq}$ comparator used for all other predicates, without requiring epoch conversion or calendar arithmetic inside the circuit.

%% ───────────────────────────────────────────────────────────────────
\subsection{Nullifiers and Holder Binding}\label{sec:nullifiers}

\paragraph{Scoped nullifiers.}
To enable double-use detection without sacrificing cross-service unlinkability, we employ \emph{scoped nullifiers}.  The nullifier is computed inside the hash sub-circuit as:
\[
  \eta \;=\; \mathsf{SHA\textnormal{-}256}\bigl(m \,\|\, c\bigr),
\]
where $m$ is the raw mdoc CBOR byte string (private input) and $c \in \{0,1\}^{64}$ is a verifier-chosen \emph{contract hash} (public input).  The contract hash serves as a domain separator --- for instance, the SHA-256 hash of the service URL or a session identifier.

This construction satisfies two properties simultaneously:
\begin{enumerate}
  \item \textbf{Replay detection within scope.}  Presenting the same credential to the same verifier (same~$c$) produces the same nullifier~$\eta$.  The verifier maintains a set of seen nullifiers and rejects duplicates.
  \item \textbf{Unlinkability across scopes.}  Presenting the same credential to two different verifiers (distinct $c_0 \neq c_1$) produces distinct nullifiers $\eta_0 \neq \eta_1$.  Since SHA-256 is modelled as a random oracle, no efficient adversary can determine whether $\eta_0$ and $\eta_1$ originate from the same credential, satisfying Definition~\ref{def:unlinkability}.
\end{enumerate}

\paragraph{Threat-model limitation.}
Because the nullifier is a deterministic function of the credential bytes~$m$ (not of a holder-held secret), an adversary who obtains~$m$ --- for instance, from a compromised device --- can compute $\eta = \mathsf{SHA\textnormal{-}256}(m \| c)$ for any scope~$c$ and check whether a published nullifier corresponds to that credential.  This is weaker than systems where nullifiers derive from a secret key held by the holder.  Incorporating a holder secret into the nullifier would require an additional private input to the circuit and a corresponding key-management protocol; we leave this extension to future work.

\paragraph{Holder binding.}
In addition to the nullifier, the circuit computes a \emph{binding hash}:
\[
  \beta \;=\; \mathsf{SHA\textnormal{-}256}\bigl(\mathsf{attr}_0[0..31]\bigr),
\]
where $\mathsf{attr}_0[0..31]$ denotes the first 32~bytes of the first requested attribute's CBOR-encoded value.  In practice, the holder controls which attribute appears first by choosing the order of the attribute request list; for cross-credential linking, the holder should place the same claim (e.g., the PID unique identifier) first in both requests.  If two credentials from different issuers happen to share the same first attribute value, the binding hashes will collide --- this is an inherent property of binding to a claim value rather than to a holder secret, and mirrors the nullifier limitation discussed above.

The binding hash serves two purposes:
\begin{enumerate}
  \item It binds the proof to a specific holder identity without revealing \emph{which} attribute was used, since $\beta$ is a one-way hash.
  \item It enables \emph{opt-in cross-credential linking}: if a holder presents two different credentials (e.g., a driving licence and a health insurance card) that share the same first attribute value, the binding hashes will match, allowing a verifier to confirm that both credentials belong to the same person --- but only if the holder chooses to present them.
\end{enumerate}

Both $\eta$ and $\beta$ are computed inside the hash sub-circuit as part of its single evaluation pass, adding only two SHA-256 invocations (each a chain of compression calls) to the existing circuit.

\paragraph{Design rationale.}
Some zero-knowledge credential systems (e.g., those based on CL-signatures~\cite{camenisch2001}) compose multiple sub-protocols via MAC chains or commitment-and-prove techniques, which increases the circuit depth or requires interactive rounds.  In \textsc{zk-eidas}, the nullifier, binding hash, and escrow digest are all computed within the \emph{same} Longfellow circuit that verifies the credential.  This single-circuit approach keeps the critical path short (depth~21 for the hash sub-circuit) and avoids the need for cross-proof consistency checks, reducing both proving time and implementation complexity.
